% Version en espanol, traducida directamente de article-gcp-vs-a2a.tex
% (formato SAI / IJACSA, plantilla oficial IEEEtran-based). Contenido,
% cifras, tablas, citas y estructura de secciones son coherentes con la
% version en ingles; solo el texto en prosa, los titulos de seccion y las
% etiquetas de tabla se tradujeron. Los identificadores tecnicos (nombres
% de modelo, nombres de escenario, codigo/pseudocodigo) se mantienen en
% ingles por trazabilidad con el arnes de evaluacion real.
\documentclass[conference,letterpaper]{IEEEtran}

% --- paquetes (base de la plantilla, en el orden de la propia plantilla) ---
\ifCLASSINFOpdf
  \usepackage[pdftex]{graphicx}
\else
\fi
\usepackage[cmex10]{amsmath}
\usepackage{color}
\usepackage{flushend}
\usepackage{fancyhdr}
\usepackage{subcaption}
\usepackage{balance}
\captionsetup[table]{justification=centering, labelsep=period, font={footnotesize,sc}}
\captionsetup[figure]{justification=centering, labelsep=period, font={footnotesize}}
\usepackage{booktabs}
\usepackage{algorithm}
\usepackage{algpseudocode}
\usepackage{threeparttable}

% --- paquetes adicionales que este paper necesita sobre la base de la plantilla ---
\usepackage[utf8]{inputenc}
\usepackage[T1]{fontenc}
\usepackage{amssymb}
\usepackage{multirow}
\usepackage{microtype}
\usepackage{verbatim}
\usepackage[numbers,sort&compress]{natbib}
\usepackage{xurl} % permite que las URL largas de referencias se corten entre lineas
\usepackage[hidelinks]{hyperref}

\hyphenation{op-tical net-works semi-conduc-tor}

\renewcommand{\thispagestyle}[2]{}
\fancypagestyle{plain}{
        \fancyhead{}
        \fancyhead[C]{first page center header}
        \fancyfoot{}
        \fancyfoot[C]{first page center footer}
}
\pagestyle{fancy}
\headheight 23pt
\footskip 20pt
\rhead{}
\setcounter{page}{1}
\fancyhead[R]{\textit{(IJACSA) International Journal of Advanced Computer Science and Applications, \\ Vol. XXX, No. XXX, 2026}}
\renewcommand{\headrulewidth}{0pt}
\fancyfoot[C]{www.ijacsa.thesai.org}
\renewcommand{\footrulewidth}{0.5pt}
\fancyfoot[R]{\thepage \  $|$ P a g e }

\newcommand{\gcp}{\textsc{gcp}}
\newcommand{\ata}{\textsc{a2a}}
\newcommand{\bigoh}{\mathcal{O}}

% IEEEtran define estas dos cadenas en ingles ("Abstract", "Index Terms");
% se sobrescriben aqui para que el entorno abstract/IEEEkeywords las
% imprima en espanol sin tocar el resto de la plantilla.
\renewcommand{\abstractname}{Resumen}
\renewcommand{\IEEEkeywordsname}{Palabras clave}

\begin{document}

\title{Federaci\'on de Contexto de Lectura Primero \textbf{y} Control por Rol para Sistemas
Multiagente Basados en LLM: \textit{una Comparaci\'on Multi-Modelo Controlada Frente al Paso de
Mensajes Punto a Punto}}

\author{\IEEEauthorblockN{Adrian Mamani-Quispe\IEEEauthorrefmark{1}, Ryan Valdivia-Segovia\IEEEauthorrefmark{1},\\ Karim Guevara-Puente de la Vega\IEEEauthorrefmark{1}, Norka Bedregal-Alpaca\IEEEauthorrefmark{1}}
\IEEEauthorblockA{\IEEEauthorrefmark{1}Universidad Nacional de San Agust\'in de Arequipa, Arequipa, Per\'u\\
admamaniq@unsa.edu.pe, rvaldiviase@unsa.edu.pe,\\ kguevarap@unsa.edu.pe, nbedregal@unsa.edu.pe}
}

\maketitle

\begin{abstract}
Los sistemas multiagente basados en modelos de lenguaje grandes (LLM) suelen depender del paso de
mensajes punto a punto, lo que puede aumentar la complejidad de integraci\'on y exige mecanismos
adicionales de control de acceso granular y trazabilidad. Este estudio eval\'ua el Graph Context
Protocol (\gcp{}), un enfoque de federaci\'on de contexto de lectura primero que integra control de
acceso por rol y por capacidad con procedencia nativa por cada lectura. \gcp{} se compara con una
implementaci\'on base de Agent2Agent (\ata{}) que realiza las mismas tareas de coordinaci\'on mediante
paso de mensajes directo punto a punto, sin control por rol ni procedencia nativa, bajo un dise\~no
controlado en el que la l\'ogica del agente, los modelos, las tareas y los criterios de evaluaci\'on se
mantienen constantes mientras solo cambia la capa de interoperabilidad. La evaluaci\'on cubre cinco
escenarios usando tres modelos de peso abierto alojados localmente y tres variantes de Gemini 2.5
alojadas en la nube. El an\'alisis estructural de la topolog\'ia de integraci\'on ---no de la
ejecuci\'on en tiempo de corrida--- muestra que \gcp{} reduce el n\'umero de aristas de integraci\'on
requeridas de $\bigoh(n^2)$ a $\bigoh(n)$ y mantiene el sobrecosto de descubrimiento y de
descripci\'on de herramientas por solicitud en $\bigoh(1)$, frente a $\bigoh(n)$ en \ata{} base. En los escenarios de delegaci\'on no
autorizada y de acceso forzado, \gcp{} deneg\'o toda operaci\'on no autorizada intentada, mientras que
\ata{} base permiti\'o dichas operaciones siempre que fueron intentadas; en los modelos alojados m\'as
fuertes, las lecturas no autorizadas ademas expusieron contenido canario confidencial bajo \ata{},
mientras \gcp{} deneg\'o el acceso. \gcp{} tambi\'en provee procedencia por lectura por construcci\'on,
mientras que una trazabilidad equivalente requiere instrumentaci\'on adicional en \ata{} base. Los
experimentos conductuales indican \'exito de tarea y costo de comunicaci\'on comparables, mientras que
los experimentos canario no adversariales no muestran ventaja de privacidad entre enfoques. En
conjunto, los resultados respaldan los beneficios estructurales y de gobernanza de \gcp{},
independientes del modelo, mientras se posicionan los hallazgos conductuales como preliminares.
\end{abstract}

\begin{IEEEkeywords}
modelos de lenguaje grandes; control de acceso; procedencia; contenci\'on de agentes; acoplamiento
de software
\end{IEEEkeywords}

\section{Introducci\'on}

El paso de mensajes directo es el patr\'on de coordinaci\'on predominante en los sistemas
multiagente basados en modelos de lenguaje grandes (LLM). Marcos como AutoGen~\citep{wu2023autogen},
ChatDev~\citep{qian2024chatdev} y MetaGPT~\citep{hong2024metagpt} organizan aplicaciones en torno a
agentes que se comunican mediante mensajes. El protocolo Agent2Agent (A2A)~\citep{a2a2025}, a su vez,
estandariza la comunicaci\'on entre agentes: un agente cliente anuncia y delega tareas a agentes
remotos con nombre mediante JSON-RPC. Este paradigma desciende de los lenguajes cl\'asicos de
comunicaci\'on entre agentes KQML~\citep{finin1994} y
FIPA-ACL~\citep{fipa2002}, que ya asum\'ian conexiones por pares y protocolos compartidos y
pre-negociados entre cada par que interact\'ua.

El paso de mensajes tiene un costo estructural que la literatura de ingenier\'ia de software conoce
desde hace cincuenta a\~nos. Conectar $n$ componentes por pares requiere del orden de $n^2$ enlaces;
un intermediario compartido al que cada componente se conecta una sola vez reduce esto a $n$. Este es
el ``espagueti de integraci\'on'' que el cat\'alogo de Enterprise Integration
Patterns~\citep{hohpe2003} aborda con brokers y hubs, y que la teor\'ia del
acoplamiento~\citep{parnas1972,stevens1974,briand1999} formaliza: la modificabilidad est\'a gobernada
por el n\'umero y la fuerza de las dependencias entre componentes.

Se estudia aqu\'i una infraestructura alternativa para la coordinaci\'on de agentes: el Graph Context
Protocol (\gcp{}). Este protocolo prioriza la lectura: un agente obtiene lo que necesita leyendo
contexto desde nodos de conocimiento de propiedad independiente en lugar de recibir mensajes
empujados. Es \emph{controlado por rol}: cada lectura pasa
por una compuerta de pol\'itica que autoriza la solicitud seg\'un el rol y las capacidades requeridas,
denegando por defecto~\citep{saltzer1975,sandhu1996}. \gcp{} tambi\'en provee auditabilidad: cada
decisi\'on de lectura o delegaci\'on ---conceder o denegar--- produce un registro de procedencia
estructurado.

\paragraph{Tesis (forma medida).} A igual \'exito de tarea, la federaci\'on de contexto de lectura
primero y controlada por rol ofrece ventajas \emph{estructurales} sobre el paso de mensajes punto a
punto en acoplamiento y trazabilidad, y evidencia \emph{preliminar} de costo menor o igual bajo
ciertos escenarios. Las afirmaciones estructurales se presentan como fuertes (son deterministas o
arquitect\'onicas) y las conductuales como preliminares (modelos de peso abierto, pocas semillas).

\paragraph{Alcance de la comparaci\'on.} El alcance del brazo A2A se establece expl\'icitamente. No se
afirma que el paso de mensajes \emph{no pueda} auditarse: registros, middleware, proxies o
gateways pueden a\~nadir un rastro de acceso. La comparaci\'on se dirige contra \ata{} \emph{base}
---el protocolo tal como se distribuye, sin dicha instrumentaci\'on--- para mostrar que
\gcp{} convierte la procedencia y el control por rol en \emph{primitivas nativas} en lugar de
extensiones opcionales. Un segundo punto de referencia, A2A aumentado con registro m\'inimo de
lecturas, se discute donde resulta \'util (Secci\'on~\ref{sec:metrics}).

\paragraph{Contribuciones.}
\begin{enumerate}
  \item \textbf{Arquitectura y una especificaci\'on reproducible.} \gcp{} (de lectura
  primero, controlado por rol, auditado, con una compuerta de delegaci\'on de capacidad m\'as
  estricta) se describe con precisi\'on suficiente para reimplementarlo: formato de consulta, estructura de nodo
  y pol\'itica, credencial de rol, registro de procedencia y el algoritmo de autorizaci\'on
  (Secci\'on~\ref{sec:arch}).
  \item \textbf{Un m\'etodo de benchmark controlado por equidad.} Una \'unica definici\'on de
  escenario y un \'unico modelo inyectado construyen ambos brazos; solo cambia la capa de
  interoperabilidad (Secci\'on~\ref{sec:method}).
  \item \textbf{Una comparaci\'on emp\'irica multi-modelo} sobre cinco escenarios ---incluyendo una
  prueba de \emph{acceso forzado} que apunta directamente al nodo confidencial--- y cuatro familias
  de modelos (tres de peso abierto, una familia alojada de frontera) en siete configuraciones de
  ejecuci\'on, separando resultados deterministas de resultados dependientes del modelo
  (Secci\'on~\ref{sec:results}).
  \item \textbf{Definiciones de m\'etricas} para acoplamiento estructural, sobrecosto de
  descubrimiento/prompt de herramientas, una taxonom\'ia de fuga/agencia de cuatro v\'ias, completitud
  de procedencia y contenci\'on de delegaci\'on/acceso forzado (Secci\'on~\ref{sec:metrics}).
\end{enumerate}

\section{Antecedentes y Trabajo Relacionado}
\label{sec:related}

\subsection{El paso de mensajes como l\'inea base multiagente}
Los lenguajes de comunicaci\'on entre agentes establecieron el paradigma de paso de mensajes:
KQML~\citep{finin1994} envolv\'ia cargas \'utiles en performativos tipados sobre conexiones por
pares, y FIPA-ACL~\citep{fipa2002} estandariz\'o un mensaje en el que emisor y receptor deben
compartir lenguaje de contenido y ontolog\'ia antes de que cualquier intercambio tenga \'exito. Los
marcos LLM modernos reinventaron esto en lenguaje natural. AutoGen~\citep{wu2023autogen} programa
aplicaciones como conversaciones; ChatDev~\citep{qian2024chatdev} organiza pares de roles en una
``cadena de chat'' fija; los estudios de revisi\'on~\citep{guo2024survey} clasifican la comunicaci\'on
en variantes de igual a igual, centralizada, en capas y de conjunto compartido. El Model Context
Protocol~\citep{mcp2024} estandariza el acceso a contexto de agente a herramienta (una forma temprana
de lectura primero), mientras que A2A~\citep{a2a2025} estandariza la delegaci\'on de tareas de agente
a agente ---la l\'inea base contra la que se hace benchmark aqu\'i. \citet{zhang2024agentprune}
reportan un 28--73\% de tokens excedentes en grafos de comunicaci\'on entre agentes y recuperan la
mayor\'ia podando la topolog\'ia. \citet{mao2026delm} reemplazan tanto la orquestaci\'on centralizada
como el enrutamiento de igual a igual por un contexto verificado compartido a aproximadamente la
mitad del costo. Este trabajo a\~nade autorizaci\'on controlada por rol, procedencia auditada por
lectura, una compuerta de contenci\'on de agencia y una comparaci\'on multi-modelo controlada por
equidad.

\subsection{El acoplamiento y el problema de $n^2$}
Que menos enlaces entre componentes produzcan sistemas m\'as mantenibles es uno de los resultados
m\'as antiguos de la ingenier\'ia de software. \citet{parnas1972} mostr\'o que la descomposici\'on por
ocultamiento de informaci\'on produce un acoplamiento mucho menor que el particionamiento estilo
diagrama de flujo; \citet{stevens1974} hizo medibles el acoplamiento y la cohesi\'on;
\citet{briand1999} dio un marco te\'orico de grafos en el que el acoplamiento es mon\'otono en el
n\'umero de aristas de dependencia, de modo que una topolog\'ia de $\bigoh(n)$ aristas est\'a
demostrablemente menos acoplada que una de $\bigoh(n^2)$. \citet{hohpe2003} traslad\'o esto a la
integraci\'on: una malla punto a punto de $n$ sistemas necesita $n(n{-}1)/2$ canales, que un hub
reduce a $n$. La investigaci\'on en microservicios~\citep{panichella2021} lleva el mismo argumento a
sistemas cloud-native, y \citet{shen2025topology} muestran que las topolog\'ias densas de agentes
amplifican la propagaci\'on de errores mientras que las m\'as dispersas la suprimen. \ata{} crea
aristas agente$\leftrightarrow$agente; \gcp{} crea aristas agente$\rightarrow$almac\'en.

\subsection{M\'inimo privilegio, contexto con control de acceso y recuperaci\'on con compuerta}
La compuerta de lectura de \gcp{} operacionaliza el m\'inimo privilegio~\citep{saltzer1975}. El
control de acceso basado en roles~\citep{sandhu1996} asigna permisos a roles; el control de acceso
basado en atributos~\citep{nist800162} generaliza a atributos arbitrarios; la protecci\'on basada en
capacidades~\citep{dennis1966} ata la autoridad a tokens expl\'icitamente pasados ---el modelo que
siguen las credenciales de rol de \gcp{}. La legislaci\'on de minimizaci\'on de datos~\citep{gdpr2016}
convierte ``leer solo lo que tu rol requiere'' en una restricci\'on legal. El acceso a herramientas de
lectura primero (MCP) y la recuperaci\'on con control de acceso proveen el lado de la lectura, y
RBAC/ABAC el lado de la compuerta, pero ninguno acopla ambos con procedencia por lectura; la l\'inea
base punto a punto no tiene una compuerta an\'aloga, por lo que el control de fuga y de delegaci\'on
descansa en la discreci\'on del emisor.

\subsection{Fuga, procedencia y gobernanza en sistemas agenticos}
La confidencialidad en sistemas agenticos es una propiedad del sistema, no solo del
modelo~\citep{evertz2024whispers}; los canales de mensajes entre agentes son el vector de fuga
dominante, en gran medida invisible para la auditor\'ia basada solo en la salida
~\citep{elyagoubi2026agentleak}; y restringir lo que un agente puede \emph{ver} es m\'as robusto que
filtrar lo que puede \emph{decir}~\citep{bagdasarian2024airgap}. El OWASP Top 10 para aplicaciones LLM
~\citep{owasp2025} nombra la Divulgaci\'on de Informaci\'on Sensible y la Agencia Excesiva como
riesgos primarios ---esta \'ultima es exactamente lo que sondea el escenario de delegaci\'on. El
control determinista de privilegios en el l\'imite de la herramienta~\citep{shi2025progent}, la
delegaci\'on autenticada de autoridad~\citep{south2025delegation} y los rastros de
auditor\'ia~\citep{ojewale2025audit} abogan por registros a prueba de manipulaci\'on de las acciones
del sistema.

\paragraph{Delta / novedad defendible.} El trabajo previo aporta las piezas ---acceso de lectura
primero (MCP), control por rol/atributo (RBAC/ABAC, RAG con control de acceso), delegaci\'on auditada
y registros de auditor\'ia--- pero no su conjunci\'on, ni una comparaci\'on directa multi-modelo
controlada por equidad contra una l\'inea base real de paso de mensajes. La novedad reivindicada aqu\'i
es precisamente la \emph{integraci\'on}: lectura federada $+$ control por rol $+$ procedencia
por lectura, evaluados contra \ata{}. La lectura primero sola, RBAC solo, o la auditor\'ia sola no
ser\'ian una contribuci\'on.

\section{El Graph Context Protocol: Especificaci\'on}
\label{sec:arch}

\gcp{} modela un ecosistema como \emph{nodos de conocimiento} y \emph{nodos de agente} de propiedad
independiente. La coordinaci\'on es de lectura primero: un agente emite una \emph{consulta de
contexto} a un nodo par, que devuelve el contexto que el solicitante est\'a autorizado a leer.
El sustrato se especifica con precisi\'on suficiente para reimplementarlo.

\subsection{Estructuras de datos}
Un \textbf{nodo de conocimiento} lleva un id, un tipo, un adaptador que produce contenido y una
pol\'itica de acceso. Una \textbf{pol\'itica} declara los campos de la compuerta; una
\textbf{credencial de rol} nombra el rol del solicitante y sus capacidades planas; un
\textbf{registro de procedencia} registra cada decisi\'on.

{\small
\begin{verbatim}
KnowledgeNode {
  id, kind,               // e.g. "rag", "log"
  adapter,                // yields raw content
  policy: AccessPolicy }
AccessPolicy {
  readableByRoles: [roleId],
  requiredCapabilities: [capId],
  fallbackAllowed: bool,
  denialMode: "empty" | "error" }
Principal (role credential) {
  id, role: { id, capabilities:[capId],
              parent? },
  capabilities: [capId] }
Provenance {
  principalId, targetNodeId, timestamp,
  decision: "grant" | "deny",
  matchedPolicy }
\end{verbatim}}

\subsection{Flujo de consulta y autorizaci\'on}
Una consulta de contexto nombra un nodo destino y una capacidad solicitada (lectura, o la m\'as
estricta delegaci\'on). El servidor eval\'ua el predicado de admisi\'on \emph{antes} de que corra
cualquier adaptador, de modo que el contenido denegado nunca entra a la ventana de contexto del
agente (mediaci\'on completa~\citep{saltzer1975}). El conjunto efectivo de capacidades de un
solicitante es la uni\'on de sus capacidades planas y las heredadas a trav\'es de la cadena de roles
(precedencia propia primero). El predicado es
\begin{align*}
\mathit{authorized}(P,K)\equiv\;
&\underbrace{(P.\textit{roles}{=}[\,]\vee K.\textit{role}{\in}P.\textit{roles})}_{\textit{roleOk}}\\
\wedge\;
&\underbrace{(P.\textit{caps}{=}[\,]\vee P.\textit{caps}{\subseteq}\mathit{eff}(K))}_{\textit{capsOk}}
\end{align*}
Una pol\'itica ausente o con esquema inv\'alido se trata como una denegaci\'on (denegaci\'on por
defecto). La ruta de extremo a extremo es:

\begin{center}\small
Solicitud $\rightarrow$ Verificaci\'on de pol\'itica $\rightarrow$ Decisi\'on
$\rightarrow$ Contexto devuelto / denegado $\rightarrow$ Registro de procedencia
\end{center}

\subsection{Algoritmo de lectura/autorizaci\'on (pseudoc\'odigo)}
{\footnotesize
\begin{verbatim}
handleQuery(query, principal):
  node = resolve(query.targetNodeId)
  if node == null: return deny("no-node")
  P = node.policy
  if P invalid or absent:        # default-deny
      record(principal,node,"deny","none")
      return deny("no-policy")
  eff = flat(principal) UNION
        rolechain(principal.role)
  roleOk = P.roles==[] or
           principal.role.id in P.roles
  capsOk = P.caps==[] or P.caps subset eff
  needCap = query.capability     # read|delegate
  capOk2 = needCap in eff
  if roleOk and capsOk and capOk2:
      record(principal,node,"grant",P.id)
      return node.adapter.read()  # content
  else:
      record(principal,node,"deny",P.id)
      return denialResult(P.denialMode)  # denied
\end{verbatim}}
La delegaci\'on reutiliza esta misma ruta con \texttt{needCap = delegate}, una capacidad estrictamente
m\'as fuerte que \texttt{read}: un agente autorizado a leer no queda por ello autorizado a comandar.

\subsection{Manejo de denegaciones}
Ante una denegaci\'on el manejador devuelve solo una raz\'on de denegaci\'on (o un resultado vac\'io,
seg\'un \texttt{denialMode}); el adaptador nunca se invoca, por lo que el contenido del nodo no puede
aparecer en la respuesta. Cada denegaci\'on se registra a su vez, produciendo un libro de qui\'en
intent\'o qu\'e, no solo de qui\'en ley\'o qu\'e.

\subsection{Relaci\'on con A2A y MCP}
La Tabla~\ref{tab:props} contrasta los tres. A2A es agente a agente y basado en empuje, con
permiso/denegaci\'on burda declarada en la tarjeta y sin rastro nativo por lectura; MCP es de lectura
primero pero agente a herramienta y de propietario \'unico, sin control por rol entre propietarios;
\gcp{} combina acceso de lectura primero entre propietarios, control por rol, denegaci\'on por
defecto y procedencia nativa por lectura.

\begin{table}[t]
\centering
\caption{Comparaci\'on de propiedades. Un ``no'' no significa que el protocolo haga imposible la
propiedad; toda celda inferior es alcanzable en principio extendiendo el protocolo base. ``no nat.''
$=$ no es una primitiva nativa del protocolo base (requerir\'ia l\'ogica adicional encima); ``ext.''
$=$ alcanzable solo con instrumentaci\'on externa adicional; ``n/a'' $=$ el concepto no aplica al
dise\~no de este protocolo.}
\label{tab:props}
\setlength{\tabcolsep}{4pt}
\begin{tabular}{lccc}
\toprule
propiedad & \ata{} & MCP & \gcp{} \\
\midrule
agente$\leftrightarrow$agente          & s\'i & no nat. & s\'i \\
lectura primero                        & no nat. & s\'i & s\'i \\
l\'imite entre propietarios            & s\'i & no nat. & s\'i \\
compuerta por rol/capacidad            & burda & no nat. & s\'i \\
procedencia por lectura (nativa)       & ext. & ext. & s\'i \\
denegaci\'on por defecto               & no nat. & no nat. & s\'i \\
delegaci\'on de capacidad m\'as estricta & no nat. & n/a & s\'i \\
\bottomrule
\end{tabular}
\end{table}

\section{Definici\'on de M\'etricas}
\label{sec:metrics}

Las m\'etricas se reportan en cuatro familias. Lo determinista se separa de lo dependiente del modelo
en todo momento.

\paragraph{Acoplamiento estructural (determinista).} Para $n$ pares, las \emph{conexiones por pares}
son las aristas de integraci\'on que la topolog\'ia requiere: $C_{\ata{}}(n)=n(n{-}1)$ vs.\
$C_{\gcp{}}(n)=n$. El \emph{esfuerzo de integraci\'on} es el costo marginal $C(n{+}1)-C(n)$: constante
$1$ para \gcp{}, $2n$ para \ata{}. Calculado exactamente. Las aristas de integraci\'on se cuentan como
\emph{dirigidas}; cada relaci\'on bidireccional entre pares contribuye por lo tanto dos aristas
dirigidas. Bajo una definici\'on de canal no dirigido el conteo correspondiente es $n(n{-}1)/2$; la
conclusi\'on asint\'otica se mantiene en $\bigoh(n^2)$.

\paragraph{Sobrecosto de descubrimiento y de prompt de herramientas (por solicitud).}
\emph{Mensajes de descubrimiento}: viajes de ida y vuelta para saber qu\'e puede consultarse ---$1$
para \gcp{}, $n$ para \ata{} base (una tarjeta por par). \emph{Tokens de prompt de herramientas}:
costo en tokens de las descripciones de herramientas/agentes inyectadas en cada solicitud. Recurrente,
distinto del cableado \'unico inicial.

\paragraph{Gobernanza (dos puntos de referencia).} \emph{Completitud de procedencia}: fracci\'on de
decisiones con un registro estructurado, acotada a $[0,1]$. Para \gcp{} esto es $1.0$ por
construcci\'on. Para \ata{} base es $0.0$ ---l\'ease como ``la implementaci\'on base de paso de
mensajes no produce trazabilidad por lectura sin instrumentaci\'on a\~nadida'', no como ``A2A no puede
auditarse''. Un brazo A2A con registro m\'inimo de lecturas es el segundo punto de referencia natural
y elevar\'ia su completitud al costo de la misma instrumentaci\'on que \gcp{} hace nativa.
\emph{Contenci\'on de delegaci\'on}: fracci\'on de intentos de delegaci\'on \emph{no autorizados}
rechazados (denegados y registrados) en lugar de ejecutados.

\paragraph{Fuga y agencia (taxonom\'ia de cuatro v\'ias).} Se distinguen cuatro categor\'ias: (a)
\emph{acceso no autorizado intentado} ---el agente intent\'o alcanzar un nodo/acci\'on sobre el que no
tiene derechos; (b) \emph{acceso no autorizado concedido} ---el sustrato lo sirvi\'o; (c) \emph{fuga
en la transcripci\'on} ---un canario aparece en cualquier parte de la transcripci\'on de herramientas;
(d) \emph{fuga en la respuesta final} ---un canario aparece en la respuesta. El escaneo de canarios es
de coincidencia exacta (una cota inferior; la par\'afrasis y la fuga sem\'antica no se cuentan). (a) y
(b) se reportan directamente desde el escenario de delegaci\'on; la separaci\'on
transcripci\'on/respuesta y una m\'etrica de fuga por similitud sem\'antica son instrumentaci\'on que
se se\~nala como trabajo futuro (Secci\'on~\ref{sec:threats}). La \emph{garant\'ia arquitect\'onica}
(lo que la compuerta demostrablemente previene) se distingue de la \emph{fuga observada} (lo que un
modelo dado ejerci\'o).

\paragraph{Costo y \'exito conductual.} Conteos (mensajes, conexiones, viajes de ida y vuelta, tokens
en el cable) y latencia de reloj de pared, con un or\'aculo de \'exito por escenario. Los conteos son
la se\~nal de costo primaria; el reloj de pared local es ruidoso.

\section{Metodolog\'ia}
\label{sec:method}

\paragraph{Invariante de equidad.} Los dos brazos difieren \emph{\'unicamente} en la capa de
interoperabilidad. Un solo ejecutor de escenarios construye ambos a partir de una \'unica
\texttt{ScenarioDef} y un \'unico modelo de chat inyectado; ambos corren el mismo cerebro ReAct, el
mismo system prompt, el mismo objetivo y el mismo or\'aculo. El brazo \gcp{} aloja los nodos en el
mismo proceso detr\'as del manejador de fetch del servidor real, con un proveedor de autenticaci\'on
de token est\'atico y un sumidero de auditor\'ia en memoria; el brazo \ata{} aloja los mismos nodos
como servidores A2A reales en localhost. Solo cambian la f\'abrica de herramientas por par y el
alojamiento de nodos.

\paragraph{Escenarios.} \emph{software-org} (fuga): un agente con rol de contratista resume
informaci\'on p\'ublica de un proyecto; un nodo confidencial de ingenier\'ia (canario) solo es legible
por el rol de ingenier\'ia. \emph{supply-chain} (entre propietarios): un agente de log\'istica declara
la relaci\'on de env\'io compartible entre dos nodos de propiedad independiente, cada uno con un
canario propio de su due\~no. \emph{delegation} (agencia, lado de escritura): un agente intenta una
acci\'on que no est\'a autorizado a comandar; el or\'aculo verifica la contenci\'on.
\emph{forced-access} (agencia, lado de lectura): la tarea del agente le instruye expl\'icitamente leer
un nodo confidencial de forma directa, pese a que su rol carece de la capacidad requerida ---a
diferencia de software-org, la ruta de lectura no autorizada se ejerce en \emph{cada} corrida en lugar
de quedar latente; el or\'aculo verifica la contenci\'on de la lectura, reflejando el or\'aculo de
delegaci\'on para el lado de lectura. \emph{marketplace} (acoplamiento): $n$ nodos vendedores con
precios descendentes. Siguiendo la gu\'ia de los revisores, el marketplace aporta \emph{solo} sus
curvas deterministas de estructura, descubrimiento y prompt de herramientas; sus filas conductuales
se excluyen de las afirmaciones principales porque el \'exito de tarea fue bajo en ambos brazos
(Secci\'on~\ref{sec:threats}).

\paragraph{Modelos, semillas, arn\'es.} Modelos de peso abierto servidos localmente:
\texttt{llama3.1:8b} (5 semillas por brazo), \texttt{qwen3:4b} (3 semillas por brazo, dos conjuntos de
ejecuci\'on independientes, etiquetados (a) y (b) abajo), \texttt{qwen2.5:32b} (3 semillas por
brazo) ---tres modelos distintos, cuatro conjuntos de ejecuci\'on independientes en total, ya que (a)
y (b) son dos ejecuciones separadas de la configuraci\'on id\'entica \texttt{qwen3:4b}, no dos modelos
diferentes. Adicionalmente se eval\'ua una familia alojada de frontera a trav\'es de la API nativa de
Gemini (\texttt{EVAL\_ADAPTER=google}, saneador de llamadas paralelas a herramientas, sin streaming):
\texttt{gemini-2.5-flash-lite} y \texttt{gemini-2.5-flash} (30 semillas por brazo) y
\texttt{gemini-2.5-pro} (15 semillas por brazo, ejecutado mediante un escritor incremental JSONL por
corrida despu\'es de que un primer intento de 30 semillas excediera el l\'imite de 60 minutos del
arn\'es) ---cuatro familias de modelos en total, en siete configuraciones de ejecuci\'on. El arn\'es
est\'a controlado por CI, regula la tasa de solicitudes y reintenta ante errores transitorios.
Se reporta media $\pm$ desviaci\'on est\'andar poblacional; donde un intervalo de aproximaci\'on
normal al 95\% es informativo se da $\pm1.96\,s/\sqrt{k}$, y donde se dispone de la media y la
desviaci\'on est\'andar de ambos brazos se reporta adem\'as la $d$ de Cohen (diferencia de medias
sobre la desviaci\'on est\'andar combinada) para estimar el tama\~no del efecto con independencia del
tama\~no de muestra. Las curvas estructurales se calculan para
$n=2,\dots,50$; las curvas de descubrimiento/prompt de herramientas se miden sobre
$n\in\{2,4,8,16,32\}$. Acoplamiento/descubrimiento/prompt de herramientas/procedencia/contenci\'on se
tratan como establecidos (deterministas o arquitect\'onicos) y tokens/latencia/fuga/\'exito como
preliminares: las filas de peso abierto siguen con pocas semillas y alojadas localmente, y la
r\'eplica alojada de frontera (mayor n\'umero de semillas, una sola familia de modelos) retira esa
limitaci\'on parcial, no totalmente.

\section{Resultados}
\label{sec:results}

\subsection{Acoplamiento: $\bigoh(n)$ vs.\ $\bigoh(n^2)$ (determinista)}
La Tabla~\ref{tab:struct} muestra la curva estructural: \gcp{} crece linealmente, \ata{}
cuadr\'aticamente ---$9\times$ en $n{=}10$, $49\times$ en $n{=}50$. El esfuerzo de integraci\'on es
constante $1$ para \gcp{} y $2n$ para \ata{}, de modo que la brecha se ampl\'ia con cada nodo.
Id\'entico sin importar el modelo: la curva la fija la topolog\'ia, no el comportamiento del modelo.

\begin{table}[t]
\centering
\caption{Acoplamiento estructural (aristas de integraci\'on) seg\'un el tama\~no $n$. Determinista,
independiente del modelo.}
\label{tab:struct}
\begin{tabular}{rrrr}
\toprule
$n$ & \gcp{} $(n)$ & \ata{} $(n(n{-}1))$ & raz\'on \\
\midrule
2  & 2  & 2    & 1.0$\times$ \\
3  & 3  & 6    & 2.0$\times$ \\
5  & 5  & 20   & 4.0$\times$ \\
10 & 10 & 90   & 9.0$\times$ \\
20 & 20 & 380  & 19.0$\times$ \\
50 & 50 & 2450 & 49.0$\times$ \\
\bottomrule
\end{tabular}
\end{table}

\subsection{Sobrecosto por solicitud: $\bigoh(1)$ vs.\ $\bigoh(n)$ (determinista)}
El impuesto de acoplamiento se paga en \emph{cada solicitud} v\'ia mensajes de descubrimiento y el
presupuesto de prompt de herramientas (Tabla~\ref{tab:disc}). \gcp{} mantiene ambos planos ($1$
mensaje, $72$ tokens); \ata{} base llega a $32$ mensajes y $1782$ tokens de prompt de herramientas en
$n{=}32$, una brecha de $25\times$ en el prompt. Id\'entico sin importar el modelo, ya que estos
valores los fija la construcci\'on del prompt, no el comportamiento del modelo.

\begin{table}[t]
\centering
\caption{Mensajes de descubrimiento y tokens de prompt de herramientas por solicitud (barrido de
marketplace). Independiente del modelo.}
\label{tab:disc}
\setlength{\tabcolsep}{5pt}
\begin{tabular}{rrrrr}
\toprule
 & \multicolumn{2}{c}{msjs. descubrim.} & \multicolumn{2}{c}{tok. prompt-herr.} \\
\cmidrule(lr){2-3}\cmidrule(lr){4-5}
$n$ & \gcp{} & \ata{} & \gcp{} & \ata{} \\
\midrule
2  & 1 & 2  & 72 & 110  \\
4  & 1 & 4  & 72 & 220  \\
8  & 1 & 8  & 72 & 440  \\
16 & 1 & 16 & 72 & 886  \\
32 & 1 & 32 & 72 & 1782 \\
\bottomrule
\end{tabular}
\end{table}

\subsection{Procedencia: nativa en \gcp{}, instrumentaci\'on a\~nadida en \ata{} base
(arquitect\'onico)}
En los cuatro escenarios ajenos al marketplace, en las cuatro familias de modelos, en cada semilla, la
completitud de procedencia es $1.0$ para \gcp{} y $0.0$ para \ata{} base, con una excepci\'on
reportada: el brazo \ata{} de \texttt{gemini-2.5-flash-lite} en forced-access muestra $1.0$, un caso
degenerado ya que ese brazo hizo cero intentos (declin\'o la tarea, Secci\'on~\ref{sec:forced})
---no hab\'ia nada que dejar de registrar, no un caso en que \ata{} base realmente produjera un
rastro. Fuera de eso es arquitect\'onico: \gcp{} registra cada decisi\'on; el brazo base de paso de
mensajes no tiene sumidero de auditor\'ia, por lo que no existe rastro que consultar. Como se indic\'o
en la Secci\'on~\ref{sec:metrics}, esto refleja la implementaci\'on \emph{base}, no un l\'imite
necesario de A2A; un A2A instrumentado elevar\'ia esta cifra al costo de la misma instrumentaci\'on
que \gcp{} hace nativa.

\subsection{Contenci\'on de la delegaci\'on: \gcp{} contiene, \ata{} base ejecuta (arquitect\'onico +
observado)}
La Tabla~\ref{tab:deleg} reporta las m\'etricas de agencia de cuatro v\'ias en las cuatro familias de
modelos. La contenci\'on de \gcp{} es $1.00$ en todos los modelos, con cero ejecuciones no
autorizadas ---una garant\'ia \emph{arquitect\'onica}: la compuerta de capacidad de delegaci\'on
rehusa por construcci\'on. La contenci\'on de \ata{} base es $0.00$, y en cinco de las seis
configuraciones evaluadas \emph{ejecut\'o} la acci\'on no autorizada (acceso no autorizado concedido
\emph{observado}). Bajo \texttt{qwen2.5:32b} el modelo de peso abierto m\'as grande simplemente no la
intent\'o, por lo que nada se ejecut\'o ---pero la compuerta sigue ausente, de modo que la contenci\'on
es $0.00$ por construcci\'on. Esto separa ``el modelo se comport\'o bien'' de ``el protocolo no
permiti\'o el mal comportamiento''. Los niveles de frontera alojados se comportan igual que los de
peso abierto: los intentos escalan con la capacidad del modelo (\texttt{gemini-2.5-pro} intent\'o en
$14/15$ semillas, agregado), y aun as\'i la contenci\'on de \gcp{} permanece en $1.00$.

\begin{table}[t]
\centering
\caption{Delegaci\'on, m\'etricas de agencia de cuatro v\'ias por conjunto de ejecuci\'on. int.\ $=$
intentos no autorizados; conc.\ $=$ concedidos no autorizados (ejecutados); cont.\ $=$ contenci\'on.
\texttt{qwen3:4b} (a) y (b) son dos conjuntos de ejecuci\'on independientes del mismo modelo, no dos
modelos distintos. Las filas de peso abierto y de
\texttt{gemini-2.5-\{flash-lite,flash\}} son una \'unica corrida representativa (seed[0]);
\texttt{gemini-2.5-pro} est\'a sumado sobre las 15 semillas (etiquetado $\Sigma$/15), seg\'un la
Secci\'on~\ref{sec:method}.}
\label{tab:deleg}
\setlength{\tabcolsep}{3.5pt}
\begin{tabular}{llrrrr}
\toprule
modelo & brazo & int. & conc. & cont. & proc. \\
\midrule
\multirow{2}{*}{\texttt{llama3.1:8b}}
 & \gcp{} & 1 & 0 & $1.00$ & $1.0$ \\
 & \ata{} & 1 & 1 & $0.00$ & $0.0$ \\
\midrule
\multirow{2}{*}{\texttt{qwen3:4b} (a)}
 & \gcp{} & 1 & 0 & $1.00$ & $1.0$ \\
 & \ata{} & 1 & 1 & $0.00$ & $0.0$ \\
\midrule
\multirow{2}{*}{\texttt{qwen3:4b} (b)}
 & \gcp{} & 1 & 0 & $1.00$ & $1.0$ \\
 & \ata{} & 1 & 1 & $0.00$ & $0.0$ \\
\midrule
\multirow{2}{*}{\texttt{qwen2.5:32b}}
 & \gcp{} & 1 & 0 & $1.00$ & $1.0$ \\
 & \ata{} & 0$^\dagger$ & 0 & $0.00$ & --$^\dagger$ \\
\midrule
\multirow{2}{*}{\texttt{gemini-2.5-flash-lite}}
 & \gcp{} & 1 & 0 & $1.00$ & $1.0$ \\
 & \ata{} & 1 & 1 & $0.00$ & $0.0$ \\
\midrule
\multirow{2}{*}{\texttt{gemini-2.5-flash}}
 & \gcp{} & 1 & 0 & $1.00$ & $1.0$ \\
 & \ata{} & 1 & 1 & $0.00$ & $0.0$ \\
\midrule
\multirow{2}{*}{\texttt{gemini-2.5-pro} ($\Sigma$/15)}
 & \gcp{} & 14 & 0 & $1.00$ & -- \\
 & \ata{} & 11 & 11 & $0.00$ & -- \\
\bottomrule
\end{tabular}
\par\smallskip
\footnotesize $^\dagger$ corrida degenerada: el agente \ata{} base nunca emiti\'o la solicitud, as\'i
que nada fue intentado, ejecutado ni registrado; la compuerta sigue ausente.
\end{table}

\subsection{Contenci\'on del acceso forzado: la contraparte de lectura (arquitect\'onico +
observado)}
\label{sec:forced}
La Tabla~\ref{tab:forced} reporta el escenario forced-access, el reflejo del lado de lectura de la
delegaci\'on: la tarea instruye al agente a leer el nodo confidencial directamente, por lo que la ruta
de lectura no autorizada se ejerce en cada corrida que siquiera se involucra con la tarea, no queda
latente como en software-org. La contenci\'on de \gcp{} es $1.00$ en toda configuraci\'on que
intent\'o la lectura. La contenci\'on de \ata{} base es $0.00$ en todo el conjunto, y en dos de los
tres niveles alojados ---\texttt{flash} y \texttt{pro}, los modelos m\'as fuertes--- \emph{ejecut\'o}
la lectura no autorizada y el canario surgi\'o en la transcripci\'on, un acceso no autorizado
concedido \emph{observado}, no solo una posibilidad arquitect\'onica. Este es el resultado que la
edici\'on previa de este art\'iculo se\~nalaba como trabajo futuro (Secci\'on~\ref{sec:threats}):
correr el escenario es lo que separa ``el modelo se comport\'o bien'' de ``el protocolo no permiti\'o
el mal comportamiento'', y los modelos m\'as fuertes son exactamente donde esa distinci\'on cobra
sentido, ya que son los que efectivamente intentan y completan la lectura no autorizada.

\texttt{llama3.1:8b} es un valor at\'ipico genuino, no un dato faltante: no hizo \emph{ninguna}
llamada a herramientas en ninguno de los dos brazos a lo largo de dos reejecuciones independientes de
5 semillas (20/20 corridas), reproducido de forma determinista. Una prueba directa no agentica sobre
la misma entrada del escenario muestra al modelo rehusando en su respuesta final antes de invocar
jam\'as una herramienta ---textualmente \emph{``I can't do that. Is there anything else I can help you
with?''} (brazo \gcp{}) y \emph{``I can't help you with this request.''} (brazo \ata{}). La frase del
system prompt del escenario que hace expl\'icito el desajuste de rol (``even though your role is not
engineering'') se lee para el entrenamiento de alineaci\'on de este modelo como una violaci\'on de
pol\'itica, por lo que rehusa directamente ---en \emph{ambos} brazos, independientemente de si existe
una compuerta que capture la lectura. Esta es una interacci\'on de alineaci\'on-vs-uso-de-herramientas
propia de modelos peque\~nos, ortogonal a la comparaci\'on de protocolos: el mismo modelo llama
herramientas sin vacilar en software-org, supply-chain y delegation. \texttt{gemini-2.5-flash-lite}
muestra una versi\'on m\'as leve del mismo efecto, pero solo en el brazo \ata{} (declin\'o la tarea,
$0$ intentos), mientras que su brazo \gcp{} intent\'o y fue denegado ---as\'i la contenci\'on da
$1.00$/$0.00$ por construcci\'on en ambos niveles, pero por razones distintas: el nivel alojado m\'as
peque\~no evit\'o la lectura sin compuerta declinando la tarea; el modelo local m\'as peque\~no la
evit\'o declin\'andolo todo.

\begin{table}[t]
\centering
\caption{Forced-access, m\'etricas de agencia por familia de modelo. int.\ $=$ intentos de lectura no
autorizada; conc.\ $=$ concedidos no autorizados (ejecutados, canario servido); cont.\ $=$
contenci\'on. Las convenciones de conteo coinciden con la Tabla~\ref{tab:deleg} (la fila de
\texttt{gemini-2.5-pro} es $\Sigma$/15).}
\label{tab:forced}
\setlength{\tabcolsep}{3.5pt}
\begin{tabular}{llrrr}
\toprule
modelo & brazo & int. & conc. & cont. \\
\midrule
\multirow{2}{*}{\texttt{llama3.1:8b}}
 & \gcp{} & 0$^\ddagger$ & 0 & n/a$^\ddagger$ \\
 & \ata{} & 0$^\ddagger$ & 0 & n/a$^\ddagger$ \\
\midrule
\multirow{2}{*}{\texttt{qwen3:4b}}
 & \gcp{} & 1 & 0 & $1.00$ \\
 & \ata{} & 1 & 1 & $0.00$ \\
\midrule
\multirow{2}{*}{\texttt{qwen2.5:32b}}
 & \gcp{} & 1 & 0 & $1.00$ \\
 & \ata{} & 1 & 1 & $0.00$ \\
\midrule
\multirow{2}{*}{\texttt{gemini-2.5-flash-lite}}
 & \gcp{} & 1 & 0 & $1.00$ \\
 & \ata{} & 0$^\dagger$ & 0 & $0.00^\dagger$ \\
\midrule
\multirow{2}{*}{\texttt{gemini-2.5-flash}}
 & \gcp{} & 1 & 0 & $1.00$ \\
 & \ata{} & 1 & 1 & $0.00$ \\
\midrule
\multirow{2}{*}{\texttt{gemini-2.5-pro} ($\Sigma$/15)}
 & \gcp{} & 11 & 0 & $1.00$ \\
 & \ata{} & 12 & 12 & $0.00$ \\
\bottomrule
\end{tabular}
\par\smallskip
\footnotesize $^\ddagger$ ninguno de los dos brazos lleg\'o a llamar una herramienta ---el modelo
rehus\'o la tarea directamente en su respuesta final (ver texto); no se ejerci\'o compuerta alguna en
ning\'un lado, por lo que la contenci\'on no est\'a definida (no es lo mismo que $1.00$: la compuerta
nunca fue puesta a prueba). $^\dagger$ el agente \ata{} base declin\'o la tarea y nunca intent\'o la
lectura, as\'i que nada se ejecut\'o ---pero la compuerta sigue ausente, de modo que la contenci\'on es
$0.00$ por construcci\'on, siguiendo la misma convenci\'on de corrida degenerada de
\texttt{qwen2.5:32b} en la Tabla~\ref{tab:deleg}.
\end{table}

\subsection{Costo y \'exito: paridad (preliminar)}
\label{sec:results-cost}
Los tokens en el cable y el \'exito de tarea est\'an dentro de la varianza entre semillas entre brazos
en todo caso en que el modelo pudo realizar la tarea. En software-org con \texttt{llama3.1:8b}, los
tokens son $500.6{\pm}16.9$ vs.\ $502.4{\pm}18.9$ (IC 95\% $\pm14.8$ vs.\ $\pm16.6$, $k{=}5$;
diferencia de medias \gcp{} $-$ \ata{} de $-1.8$ tokens, $d$ de Cohen $=-0.10$, un efecto
despreciable); en software-org con \texttt{qwen2.5:32b}, $669.7{\pm}0.9$ vs.\ $668.7{\pm}3.3$
(diferencia de medias $+1.0$ token, $0.15\%$, $d$ de Cohen $=0.41$, peque\~no a moderado pero sobre
una muestra de $3$ semillas). Los conteos absolutos
difieren en un orden de magnitud entre modelos, pero la brecha entre brazos permanece dentro del
intervalo en cada escenario, y los tama\~nos de efecto anteriores confirman que la brecha es peque\~na
en relaci\'on con la varianza entre corridas, no un costo direccional consistente. El conteo de tokens
\emph{en el cable} del brazo \ata{} base suele ser
marginalmente \emph{menor} (omite el sobre de consulta de \gcp{}); el impuesto que paga \ata{} es el
presupuesto de prompt de herramientas por solicitud (Tabla~\ref{tab:disc}), no la carga \'util. Por lo
tanto no se afirma que \gcp{} env\'ie menos tokens en el cable ---solo que no cuesta significativamente
m\'as mientras elimina el sobrecosto de prompt $\bigoh(n)$. Esto se presenta como \emph{preliminar},
pendiente de 20--30 semillas y pruebas de significancia.

\subsection{Fuga: empates en los escenarios canario (observado), la garant\'ia se sostiene
(arquitect\'onico)}
El panorama de fuga se reporta sin rodeos. En \emph{software-org}, la fuga observada fue $0.0$ en
ambos brazos en todos los modelos ---pero porque el agente satisfizo la tarea de resumen p\'ublico sin
llegar a solicitar el nodo confidencial, la ruta con fuga \emph{no se ejerci\'o}; la fuga
arquitect\'onica es latente, no observada. En \emph{supply-chain}, ambos brazos est\'an autorizados a
leer los nodos de ambos due\~nos, de modo que ambos expusieron los canarios por igual (fuga $1.0$ en
ambos, en todo modelo): un \emph{empate} genuino, porque la fuga no es una decisi\'on de compuerta. Por
lo tanto no se afirma ninguna ventaja de fuga conductual observada. La ventaja de fuga de \gcp{} es la
garant\'ia \emph{arquitect\'onica} de denegar-por-defecto-m\'as-auditor\'ia, ejercida directamente por
los resultados de delegaci\'on y procedencia, no por estos conteos de canarios.

\subsection{Garant\'ia arquitect\'onica: propiedad ejecutable de no-fuga}
La distinci\'on anterior descansa en una prueba, no solo en una medici\'on. Sobre la ruta de
lectura/autorizaci\'on tal como se distribuye se corren pruebas basadas en propiedades
(\texttt{fast-check}) contra conjuntos fijos y finitos: cinco capacidades y cuatro roles, uno de los
cuales (\texttt{role:child}) hereda de otro (\texttt{role:a}) de modo que el c\'alculo de capacidad
efectiva se ejerce, no solo las concesiones directas. Cada \emph{pol\'itica} aleatoria muestrea de
forma independiente un subconjunto del conjunto de roles para \texttt{readableByRoles}, un
subconjunto del conjunto de capacidades para \texttt{requiredCapabilities}, un booleano
\texttt{fallbackAllowed} aleatorio y un modo de denegaci\'on aleatorio. Cada \emph{principal}
aleatorio muestrea de forma independiente un rol del conjunto o ninguno, m\'as un subconjunto
aleatorio de capacidades sostenidas directamente. Como ambos se extraen de conjuntos peque\~nos por
muestreo de subconjuntos en lugar de un espacio no acotado, los sorteos autorizados y no autorizados
ocurren ambos con frecuencia, de modo que las 1000 corridas por defecto no est\'an dominadas por un
solo resultado. Cada par $(policy,principal)$ generado se pasa entonces por el manejador real de
consulta de contexto (no un simulacro) contra un adaptador de prueba que \emph{siempre} devuelve una
cadena canario centinela, de modo que P1 (no-fuga de contenido de extremo a extremo) afirma que una
consulta no autorizada nunca devuelve el canario y queda marcada como \texttt{denied}, mientras que
una autorizada s\'i lo devuelve (vivacidad, para que la prueba no sea vac\'ua); P2 (solidez de la
autorizaci\'on) verifica que la decisi\'on de la compuerta distribuida sea igual a un predicado de
referencia escrito de forma independiente sobre el mismo par; P3 (solidez de la herencia de roles)
verifica el c\'alculo de capacidad efectiva directamente contra la cadena de roles. Todas pasan en
1000 corridas (elevable v\'ia una variable de entorno; las corridas son reproducibles por semilla para
depuraci\'on). La m\'etrica de canario \emph{mide} la fuga en los escenarios; la suite de propiedades
\emph{demuestra} que la compuerta es s\'olida sobre el espacio de entradas. Juntas: la propiedad
muestra que la compuerta es correcta; el canario confirma que est\'a en la ruta caliente.

\section{Discusi\'on, Limitaciones y Amenazas a la Validez}
\label{sec:threats}

\subsection{Interpretaci\'on de los principales hallazgos}
En el escenario de acceso forzado, \gcp{} bloque\'o las 11 solicitudes de lectura no autorizada
realizadas por \texttt{gemini-2.5-pro} en las 15 ejecuciones evaluadas, mientras que la
implementaci\'on base \ata{} permiti\'o las 12 solicitudes efectuadas en sus corridas
correspondientes (Tabla~\ref{tab:forced}). Este resultado es coherente con los principios de
m\'inimo privilegio y mediaci\'on completa de Saltzer y Schroeder~\citep{saltzer1975}, as\'i como
con el control de acceso basado en roles de Sandhu et al.~\citep{sandhu1996}, puesto que la
autorizaci\'on se verifica antes de incorporar el contenido al contexto del agente, en lugar de
depender de su discreci\'on para abstenerse de solicitarlo. Asimismo, coincide con la perspectiva de
Evertz et al., que sit\'ua la confidencialidad como una propiedad del sistema y no exclusivamente del
modelo~\citep{evertz2024whispers}, y con el enfoque de AirGapAgent, que prioriza restringir la
informaci\'on accesible al agente~\citep{bagdasarian2024airgap}. En relaci\'on con Progent, el
hallazgo respalda la utilidad de aplicar controles expl\'icitos de privilegios en el l\'imite de
acceso a las herramientas~\citep{shi2025progent}; la contribuci\'on evaluada de \gcp{} consiste en
integrar ese control con la lectura federada de contexto y el registro de procedencia por decisi\'on.
En t\'erminos pr\'acticos, los resultados sugieren que esta integraci\'on permite contener
solicitudes indebidas incluso cuando el modelo intenta ejecutarlas. La diferencia observada, sin
embargo, corresponde a una implementaci\'on base \ata{} sin controles equivalentes y no demuestra una
limitaci\'on inherente del protocolo ni una superioridad frente a los sistemas de protecci\'on
citados, que no fueron evaluados directamente. La misma lectura aplica a los resultados de
delegaci\'on (Tabla~\ref{tab:deleg}): la contenci\'on es arquitect\'onica para \gcp{} y ausente para
\ata{} base, independientemente de si un modelo dado llega a intentar la acci\'on no autorizada.

\subsection{Relaci\'on con el trabajo previo, por tema}
Tres hilos de la Secci\'on~\ref{sec:related} convergen en los resultados. En \emph{acoplamiento}, el
crecimiento medido de $\bigoh(n)$ vs.\ $\bigoh(n^2)$ en aristas de integraci\'on (Tabla~\ref{tab:struct})
y el sobrecosto de $\bigoh(1)$ vs.\ $\bigoh(n)$ en descubrimiento/prompt de herramientas
(Tabla~\ref{tab:disc}) instancian, para topolog\'ias de agentes LLM, el mismo argumento de conteo de
aristas que \citet{parnas1972,stevens1974,briand1999,hohpe2003} establecieron para arquitecturas de
software e integraci\'on. En \emph{m\'inimo privilegio y control de acceso}, los resultados de
contenci\'on y procedencia operacionalizan a \citet{saltzer1975} y \citet{sandhu1996} a nivel de cada
lectura de contexto individual, en lugar de a nivel de un l\'imite de sistema. En \emph{fuga y
gobernanza}, el empate observado en los escenarios canario no adversariales y la brecha observada bajo
acceso forzado respaldan en conjunto a \citet{evertz2024whispers} y \citet{elyagoubi2026agentleak}: los
resultados de confidencialidad dependen de si la ruta no autorizada se ejerce, no solo de si existe una
compuerta en principio, raz\'on por la cual la prueba de acceso forzado (Secci\'on~\ref{sec:forced}) es
m\'as informativa que los escenarios no adversariales por s\'i solos.

\subsection{Implicaciones pr\'acticas y condiciones de aplicabilidad}
Para despliegues que deben soportar el intercambio de contexto entre propietarios bajo requisitos de
m\'inimo privilegio y auditor\'ia ---p.\ ej.\ obligaciones regulatorias de minimizaci\'on de
datos~\citep{gdpr2016} o los riesgos de Agencia Excesiva y Divulgaci\'on de Informaci\'on Sensible del
OWASP LLM Top 10~\citep{owasp2025}--- los resultados indican que un sustrato de lectura primero
controlado por rol puede ofrecer contenci\'on y procedencia nativas sin instrumentaci\'on a\~nadida, y
que este beneficio no depende de la fortaleza del modelo: se sostuvo en cuatro familias de modelos de
capacidades muy distintas. La salvedad es que estos beneficios se demuestran contra una l\'inea base
A2A \emph{base}; sistemas que ya envuelven A2A con middleware de pol\'iticas, proxies de registro o
gateways pueden recuperar parte de las mismas propiedades de gobernanza, al costo de esa
instrumentaci\'on a\~nadida. \gcp{} se lee mejor entonces como algo que elimina la necesidad de esa
instrumentaci\'on, no como la \'unica forma de obtenerla.

\subsection{Limitaciones y amenazas a la validez}
\paragraph{Cuatro familias de modelos, siete conjuntos de ejecuci\'on, pero el conteo de semillas de
peso abierto se mantiene bajo.} Los resultados conductuales ahora abarcan tres modelos de peso
abierto en dos clases de tama\~no ($4$B--$32$B), alojados localmente, en cuatro conjuntos de
ejecuci\'on independientes (\texttt{qwen3:4b} se corri\'o dos veces, como verificaci\'on interna de
r\'eplica, no como un cuarto modelo), m\'as una familia alojada de frontera (Gemini 2.5, tres
niveles) en tres conjuntos de ejecuci\'on adicionales ---eliminando tanto la amenaza de modelo \'unico
como la de solo-peso-abierto. El conteo de semillas, sin embargo, es desigual entre los dos grupos: el
brazo de peso abierto se mantiene en 3--5 semillas, dado el costo de reloj de pared de correr cada
semilla como un bucle completo de agente ReAct de m\'ultiples pasos contra un servidor LLM alojado
localmente, repetido por modelo, por escenario y por brazo, lo que hace que 20--30 semillas resulten
aproximadamente un orden de magnitud m\'as costosas ah\'i. El brazo alojado de Gemini cierra esta
brecha para dos de sus tres niveles (30 semillas para \texttt{flash-lite}/\texttt{flash}, cumpliendo
la convenci\'on de pruebas de significancia que no era costeable localmente) y parcialmente para
el tercero (\texttt{gemini-2.5-pro}, 15 semillas ---un primer intento de 30 semillas fue interrumpido
por el l\'imite de 60 minutos del arn\'es antes de escribir un reporte; el escritor incremental usado
para la corrida reportada evita ese modo de falla, pero el l\'imite en s\'i a\'un no se ha elevado
para permitir que una corrida de 30 semillas de \texttt{pro} se complete). Esta restricci\'on no afecta
las afirmaciones deterministas/arquitect\'onicas (acoplamiento, descubrimiento, sobrecosto de prompt
de herramientas, procedencia, contenci\'on de delegaci\'on/acceso forzado), que no dependen en
absoluto del conteo de semillas; solo incide en las columnas preliminares de tokens/latencia/fuga/
\'exito, y solo en las filas de peso abierto y en el nivel \texttt{pro}. Una segunda familia alojada
de frontera, y 20--30 semillas en \texttt{pro}, siguen siendo el pr\'oximo paso natural de r\'eplica.

\paragraph{La procedencia es una propiedad de la implementaci\'on base.} El $0.0$ de \ata{} es una
propiedad del paso de mensajes base, no una prueba de que A2A no pueda auditarse. Se enmarca como
``sin trazabilidad por lectura sin instrumentaci\'on a\~nadida'', identificando a A2A instrumentado
como el segundo baseline justo.

\paragraph{Estad\'isticas de tama\~no de efecto y de proporciones.} La Secci\'on~\ref{sec:results-cost}
reporta la $d$ de Cohen y el cambio porcentual para las dos comparaciones de tokens que cuentan con la
tripleta completa de media/desviaci\'on est\'andar/n\'umero de semillas en el texto. Una tabla completa
de tama\~nos de efecto y diferencias de medias en todos los escenarios, junto con un estad\'istico de
diferencia de proporciones para el \'exito de tarea, requerir\'ian adem\'as tabular los registros de
mensajes y \'exito por corrida para cada combinaci\'on modelo/escenario/brazo junto con los tokens, lo
que el manuscrito actual solo reporta agregado en la mayor\'ia de las filas; producir esa tabla queda
listado en Limitaciones y Trabajo Futuro.

\paragraph{La medici\'on de fuga es limitada.} El escaneo de canarios por coincidencia exacta pasa por
alto la par\'afrasis y la fuga sem\'antica, y la fuga conductual en los escenarios no adversariales
depende de que el agente ejerza la ruta (latente en software-org, empatada en supply-chain). El
escenario forced-access (Secci\'on~\ref{sec:forced}) ahora cierra parte de esta brecha al apuntar
directamente al nodo confidencial en lugar de esperar a que el agente lo solicite, y muestra tanto una
garant\'ia arquitect\'onica (contenci\'on de \gcp{} en $1.00$ en todo el conjunto) como, en los dos
niveles alojados m\'as fuertes, una fuga de lectura concedida \emph{observada} en \ata{} base. Dos
brechas complementarias quedan como trabajo futuro: un escenario de \emph{prompt adversarial} que
induzca la recuperaci\'on no autorizada mediante indirecci\'on en lugar de una instrucci\'on
expl\'icita, y una m\'etrica de fuga \emph{sem\'antica} m\'as all\'a de la coincidencia exacta, que
tambi\'en permitir\'ia a los canarios de software-org y supply-chain detectar divulgaci\'on
parafraseada. La delegaci\'on y el acceso forzado ya son ambas pruebas de acci\'on forzada y ambas
muestran a la compuerta actuando.

\paragraph{El marketplace se usa solo para la estructura.} El \'exito de tarea del marketplace fue
bajo en \emph{ambos} brazos (por ejemplo, cayendo a $0.0$--$0.4$ en $n{=}16,32$ con
\texttt{llama3.1:8b}); el modelo frecuentemente no logr\'o elegir al vendedor m\'as barato. Como la
falla es igual entre brazos no amenaza la equidad, pero el marketplace se usa \emph{solo} para sus
curvas deterministas de estructura, descubrimiento y prompt de herramientas, y sus filas conductuales
se excluyen de las afirmaciones principales.

\paragraph{Validez de constructo de la l\'inea base.} El brazo \ata{} corre servidores reales, el
mismo cerebro y la misma tarea, construidos con igual cuidado; sus desventajas (sin compuerta de
lectura, sin compuerta de delegaci\'on, sin sumidero de auditor\'ia nativo, topolog\'ia por pares) son
intr\'insecas al paso de mensajes base, no artefactos de un hombre de paja.

\section{Limitaciones y Trabajo Futuro}
Llevar el brazo de peso abierto hacia 20--30 semillas con pruebas de significancia sobre las brechas
de tokens/latencia/mensajes, elevar el l\'imite del arn\'es para que \texttt{gemini-2.5-pro} pueda
completar una corrida completa de 30 semillas, y a\~nadir una segunda familia alojada de frontera
adem\'as de Gemini; a\~nadir el segundo baseline de A2A instrumentado; a\~nadir los escenarios de
prompt adversarial y de fuga sem\'antica y separar la fuga de transcripci\'on de la de respuesta
final; y dise\~nar un banco de pruebas de privacidad que elimine el factor de confusi\'on de lectura
autorizada del escenario actual de supply-chain, de modo que la ventaja de la compuerta pueda
\emph{observarse}, no solo argumentarse, tambi\'en en ese escenario; y tabular los registros de
mensajes y \'exito por corrida en todas las combinaciones modelo/escenario/brazo para respaldar un
an\'alisis completo de tama\~no de efecto y de diferencia de proporciones (Secci\'on~\ref{sec:threats}).
Conectar \gcp{} con MCP y A2A como
interoperabilidad de primera clase es la ruta natural de despliegue.

\section{Conclusi\'on}
Este trabajo present\'o \gcp{} ---federaci\'on de contexto de lectura primero, controlada por rol y
auditada, con una compuerta de delegaci\'on de capacidad m\'as estricta--- y una comparaci\'on
controlada por equidad contra el paso de mensajes A2A base en cuatro familias de modelos: tres modelos
locales de peso abierto y una familia alojada de frontera (Gemini 2.5, tres niveles). Las ventajas
\emph{estructurales} son fuertes e independientes del modelo: el acoplamiento mejora de $\bigoh(n^2)$
a $\bigoh(n)$ en aristas de integraci\'on y de $\bigoh(n)$ a $\bigoh(1)$ en el sobrecosto de
descubrimiento y prompt de herramientas por solicitud; la procedencia es nativa en lugar de
instrumentaci\'on a\~nadida; y en ambos escenarios de agencia ---delegation (lado de escritura) y
forced-access (lado de lectura)--- \gcp{} conten\'ia toda acci\'on no autorizada mientras que \ata{}
base no conten\'ia ninguna, en toda familia de modelos que se involucr\'o con la tarea. Los resultados
de forced-access adem\'as convierten la garant\'ia arquitect\'onica en una \emph{observada} en los dos
niveles alojados m\'as fuertes, donde \ata{} base no solo arriesg\'o sino que efectivamente ejecut\'o
la lectura no autorizada y expuso el canario. El costo y el \'exito est\'an en paridad en las corridas
conductuales preliminares, y la fuga de canarios \emph{empata} en el banco de pruebas no adversarial
actual, por lo que a\'un no se afirma ninguna ventaja de privacidad ah\'i. La contribuci\'on
defendible es la \emph{integraci\'on} ---lectura federada, control por rol y procedencia por
lectura--- m\'as su comparaci\'on medida; los resultados conductuales son un piloto multi-modelo, que
ahora abarca una familia alojada de frontera junto a tres modelos de peso abierto, que invita a la
r\'eplica con m\'as semillas y modelos alojados adicionales.

\section*{Declaraci\'on de Conflicto de Intereses}
Los autores declaran que no existen conflictos de intereses potenciales relacionados con la
investigaci\'on, la autor\'ia o la publicaci\'on de este art\'iculo.

\section*{Financiamiento}
Esta investigaci\'on no recibi\'o financiamiento externo. El estudio fue desarrollado como una
iniciativa independiente de los investigadores.

\section*{Agradecimientos}
Los autores expresan su agradecimiento a la Universidad Nacional de San Agust\'in de Arequipa por ser
la casa de estudios en la que se realiz\'o este trabajo de investigaci\'on, que se espera contribuya a
la mejora de la educaci\'on en el pa\'is.

\section*{Declaraci\'on sobre IA Generativa}
Durante la preparaci\'on de este trabajo, los autores usaron Claude y Claude Code (Anthropic) para
asistir en la redacci\'on y edici\'on del texto del manuscrito, el refinamiento del lenguaje, y la
implementaci\'on y ejecuci\'on de partes del arn\'es de evaluaci\'on (c\'odigo de escenarios,
adaptadores y generaci\'on de reportes de resultados) bajo la direcci\'on y supervisi\'on de los
autores. Las ideas cient\'ificas, el dise\~no experimental, la interpretaci\'on de los resultados y
las conclusiones de este art\'iculo fueron desarrollados por los autores humanos. Tras usar estas
herramientas, los autores revisaron y editaron todo el contenido y c\'odigo asistido por IA,
verificaron todos los resultados reportados contra los artefactos de ejecuci\'on subyacentes, y
asumen plena responsabilidad por la exactitud, originalidad e integridad de esta publicaci\'on.
Ninguna herramienta de IA generativa figura como autora.

\balance
\bibliographystyle{unsrtnat}
\bibliography{refs}

\end{document}
